%\documentclass{beamer}
%\usepacakge[UTF8]{ctex}
%----------------------------------------------------------------------------------------
%    PACKAGES AND THEMES
%----------------------------------------------------------------------------------------

\documentclass[aspectratio=169,xcolor=dvipsnames]{beamer}
\usepackage[UTF8]{ctex}
\usetheme{SimplePlus}
\setbeamertemplate{footline}{}

\usepackage{hyperref}
\usepackage{graphicx} % Allows including images
\usepackage{booktabs} % Allows the use of \toprule, \midrule and \bottomrule in tables

\newcommand\Ty{\mathtt{Ty}}
\newcommand\Ter{\mathtt{Ter}}
\newcommand\ttp{\mathtt{p}}
\newcommand\ttq{\mathtt{q}}
\newcommand\bfone{\mathbf{1}}
\newcommand\Set{\mathbf{Set}}
\newcommand\C{\mathbf{C}}
\newcommand\extop{{\scalebox{0.4}[1]{\_}}.{\scalebox{0.4}[1]{\_}}}

%----------------------------------------------------------------------------------------
%    TITLE PAGE
%----------------------------------------------------------------------------------------

\title{族范畴}
%\subtitle{Subtitle}

%\author{Pin-Yen Huang}

%\institute
%{
%    Department of Computer Science and Information Engineering \\
%    National Taiwan University % Your institution for the title page
%}
%\date{\today} % Date, can be changed to a custom date
\date{}

%----------------------------------------------------------------------------------------
%    PRESENTATION SLIDES
%----------------------------------------------------------------------------------------

\begin{document}

\begin{frame}
    % Print the title page as the first slide
    \titlepage
\end{frame}

\begin{frame}{元素范畴}
  \begin{itemize}
    \item $P : \C^{op} \to \Set$ 是一个预层 \pause
    \item $P$ 的元素范畴 $\int P$ \pause，对象为 $(X,p)$\pause(其中 $X\in\C$, $p\in P(X)$); \pause 态射 $f : (X,p)\to(X',p')$\pause($f$是$\C$中的态射且满足$P(f)(p')=p$)
  \end{itemize}
\end{frame}

\begin{frame}{族范畴}
\begin{itemize}
  \item 上下文范畴$\C$\pause(对象代表类型论的context,态射代表substitution). \pause 这个范畴要有一个终对象$\bfone$, 表示空的上下文.
  \item 函子$\Ty:\C^{op}\to\Set$\pause (对上下文$\Gamma$, $\Ty(\Gamma)$代表上下文$\Gamma$中的类型).\pause 对态射$\gamma:\Delta\to\Gamma$, 类型$A\in\Ty(\Gamma)$, 用 $A[\gamma]$ 表示 $\Ty(\gamma)(A)\in\Ty(\Delta)$.\pause
  \item 函子$\Ter:(\int\Ty)^{op}\to\Set$ (对每个上下文 $\Gamma$, 类型$A\in\Ty(\Gamma)$, $\Ter(\Gamma,A)$表示在上下文$\Gamma$中类型为$A$的项).\pause ($\Ter(\Gamma,A)$也记为$\Ter(\Gamma\vdash A)$) \pause 取$\int\Ty$ 中的态射 $\gamma:(\Delta,B)\to(\Gamma,A)$ \pause(那么, $\gamma:\Delta\to\Gamma$是$\C$ 中的态射, $B=A[\gamma]$). \pause $\Ter(\gamma):\Ter(\Gamma,A)\to\Ter(\Delta,B)$. 那么对$a\in\Ter(\Gamma\vdash A)$, 用$a[\gamma]$表示$\Ter(\gamma)(a)\in\Ter(\Delta\vdash A[\gamma])$.\pause
  \item 上下文扩展操作$\extop$ \pause 对$\Gamma\in\C$, $A\in\Ty(\Gamma)$, 返回对象 $\Gamma.A\in\C$, 并且有一个态射$\ttp:\Gamma.A\to\Gamma$ \pause (表示context weakening), \pause 一个项$\ttq\in\Ter(\Gamma.A\vdash A[\ttp])$ \pause (表示新添加的类型为$A$的变量). \pause 这个操作满足以下万有性质: 对任意替换$\gamma:\Delta\to\Gamma$, 项 $a\in\Ter(\Delta\vdash A[\gamma])$, 存在唯一态射$\langle\gamma,a\rangle:\Delta\to\Gamma.A$, 使得 $\ttp\circ\langle\gamma,a\rangle=\gamma$, $\ttq[\langle\gamma,a\rangle]=a$ \pause (此处两个$\langle\gamma,a\rangle$意义不一样, $\ttp\circ\langle\gamma,a\rangle=\gamma$ 中的 $\langle\gamma,a\rangle$表示$\C$中态射, $\ttq[\langle\gamma,a\rangle]=a$ 中的 $\langle\gamma,a\rangle$ 表示 $\int\Ty$ 中的态射 $\langle\gamma,a\rangle : (\Delta,B)\to (\Gamma.A,A[\ttp])$, 其中 $B=A[\ttp][\langle\gamma,a\rangle]=A[\ttp\circ\langle\gamma,a\rangle]=A[\gamma]$)
\end{itemize}
\end{frame}

\end{document}